Dựa trên số liệu thu được từ Hình \ref{fig:quy_dao} và Bảng \ref{tab:anh_huong_luc_can}, có thể rút ra các nhận xét sau:

\begin{itemize}
\item \textbf{Sự suy giảm tầm xa và độ cao:} Khi xuất hiện lực cản không khí, cả tầm xa ($R$) và độ cao cực đại ($H_{\max}$) đều giảm mạnh so với trường hợp lý tưởng. Cụ thể, với $n=1$ và góc $45^\circ$, tầm xa giảm từ 254,82 m xuống còn 47,88 m (giảm hơn 80\%).
    
\item \textbf{Vai trò của bậc lũy thừa $n$:} Khi $n$ tăng (từ 1 đến 2), lực cản tác động lên vật thể càng lớn, đặc biệt ở giai đoạn đầu của quỹ đạo khi vận tốc còn cao. Điều này khiến vật thể mất động năng nhanh chóng, dẫn đến quỹ đạo bị “ép dẹt” và thời gian bay ($T$) rút ngắn đáng kể.
    
\item \textbf{Sự thay đổi góc bắn tối ưu:} Trong môi trường không có lực cản, góc bắn $45^\circ$ cho tầm xa lớn nhất. Tuy nhiên, khi có lực cản (đặc biệt khi $n$ lớn), góc bắn tối ưu có xu hướng giảm xuống dưới $45^\circ$.
    
\item \textbf{Hình dạng quỹ đạo:} Khác với đường parabol đối xứng trong chân không, quỹ đạo khi có lực cản trở nên bất đối xứng. Đoạn đi xuống dốc hơn đoạn đi lên do vận tốc theo phương ngang bị suy giảm liên tục bởi lực cản.
\end{itemize}
